\documentclass{article}

% Use [accepted]{icml2025} for the camera-ready (non-anonymous) version,
% \usepackage{icml2025} for the anonymous submission version.
\usepackage[accepted]{icml2025}

\usepackage{microtype}
\usepackage{graphicx}
\usepackage{booktabs}
\usepackage{amsmath}
\usepackage{amssymb}
\usepackage{algorithm}
\usepackage{algorithmic}
\usepackage{hyperref}
\usepackage{url}

\icmltitlerunning{Hardware-Aware Adaptive LoRA Rank Allocation}

\begin{document}

\twocolumn[
\icmltitle{Hardware-Aware Adaptive LoRA Rank Allocation:\\
           Trading Importance for Cost in Parameter-Efficient Fine-Tuning}

\icmlsetsymbol{equal}{*}

\begin{icmlauthorlist}
\icmlauthor{David Wang}{cu}
\end{icmlauthorlist}

\icmlaffiliation{cu}{Course project, Advanced Machine Learning Systems}

\icmlcorrespondingauthor{David Wang}{notchobrinewr@gmail.com}

\icmlkeywords{LoRA, parameter-efficient fine-tuning, adaptive rank,
              hardware-aware training, systems for ML}

\vskip 0.3in
]

\printAffiliationsAndNotice{}

\begin{abstract}
Low-rank adaptation (LoRA) reduces the number of trainable parameters
needed to fine-tune a transformer, but once the total rank budget is
fixed, \emph{how} that budget is distributed across modules becomes its
own decision. Standard LoRA splits it uniformly; AdaLoRA reallocates it
during training by SVD-pruning singular values. Neither accounts for the
fact that the same unit of rank costs different numbers of trainable
parameters at different modules. We study a simple allocator that mixes
gradient-based importance with explicit per-rank parameter cost,
$s_i = g_i / c_i^{\alpha}$, applied once at a warmup-to-stage-2 transition
and held fixed thereafter. On SST-2 with DistilBERT under a fixed
$192$-rank budget, the allocator reaches the lowest trainable-parameter
count of any method evaluated ($1{,}081{,}346$ vs.\ AdaLoRA's
$1{,}366{,}562$), comparable accuracy with the lowest seed-to-seed
variance, and a $\sim$$19\times$ reduction in scheduler overhead relative
to AdaLoRA. An $\alpha$-sweep ($0.0 \to 0.5 \to 1.0$) shows attention rank
share growing monotonically from $39.1\%$ to $56.1\%$ at constant
gradient signal, providing direct evidence that the cost term reshapes
allocation independent of the importance proxy. The picture against
Uniform LoRA is mixed --- peak memory and time-to-target favor Uniform
on this task --- which we report rather than retune.
\end{abstract}

\section{Introduction}
\label{sec:intro}

Low-Rank Adaptation (LoRA) \cite{hu2022lora} fine-tunes a frozen
pretrained transformer by inserting trainable rank-$r$ adapters
$\Delta W = BA$ into selected linear layers. Once the total rank budget
$B = \sum_i r_i$ is fixed, the allocation of $B$ across modules is itself
a design choice. Standard LoRA uses uniform $r_i$. AdaLoRA
\cite{zhang2023adalora} reallocates rank during training by pruning
SVD components that contribute least to the loss. Both view rank as
fungible across modules.

It is not. On DistilBERT, adding rank to a 768$\to$3072 feed-forward
projection produces $2.5\times$ as many trainable parameters as adding
rank to a 768$\to$768 attention projection. A unit of rank that costs
3840 parameters at \texttt{ffn.lin1} and 1536 at \texttt{attn.q\_lin}
should arguably not be priced the same when the budget is parameter-
or memory-bound.

This paper studies a small modification: divide the importance signal by
hardware cost. For each LoRA target module $i$,
\begin{equation}
s_i \;=\; \frac{g_i}{c_i^{\alpha}},
\label{eq:rule}
\end{equation}
where $g_i$ is an exponential moving average of LoRA gradient norms
collected during a warmup phase, $c_i = d_i^{\text{in}} + d_i^{\text{out}}$
is the per-rank parameter cost, and $\alpha \in [0,1]$ controls how
strongly cost is penalized. Rank is allocated proportionally to $s_i$
under a fixed total budget. Setting $\alpha{=}0$ recovers a gradient-only
ablation; $\alpha{=}1$ is the headline hardware-aware variant.

We report the systems-level consequences of this rule on SST-2
\cite{socher2013sst} with DistilBERT \cite{sanh2019distilbert} against
Uniform LoRA, AdaLoRA \cite{zhang2023adalora} (PEFT
implementation~\cite{mangrulkar2022peft}), and the $\alpha{=}0$
gradient-only ablation. Our contributions are:
\begin{itemize}
\item the allocation rule of Eq.~\ref{eq:rule} together with a
      budget-preserving discrete allocator (\S\ref{sec:method});
\item a systems-level evaluation logging peak memory, throughput,
      wall-clock, time-to-target, and \emph{scheduler overhead} ---
      the latter being the project's clearest systems finding
      (\S\ref{sec:results-systems});
\item an $\alpha$-sweep that isolates the cost term, showing it
      monotonically reshapes attention vs.\ FFN rank share at constant
      accuracy and constant gradient signal
      (\S\ref{sec:results-alpha}).
\end{itemize}
We make no claim of beating AdaLoRA on accuracy. The point is to
characterise the practical-efficiency picture under a constrained budget.

\section{Related Work}
\label{sec:related}

\textbf{LoRA} \cite{hu2022lora} introduced low-rank adapter
fine-tuning with uniform rank. \textbf{AdaLoRA} \cite{zhang2023adalora}
adaptively budgets rank across modules using an SVD-based importance
score, pruning unimportant components throughout training. Our method
differs in two ways: the importance signal is a Frobenius-norm gradient
EMA rather than an SVD criterion, and we explicitly divide that signal
by per-rank parameter cost. \textbf{LoRA+} \cite{hayou2024loraplus}
addresses a different axis (separate learning rates for the $A$ and $B$
factors) and is orthogonal to budget allocation. \textbf{QLoRA}
\cite{dettmers2023qlora} composes LoRA with 4-bit quantisation of the
base model, again orthogonal to rank distribution. \textbf{GaLore}
\cite{zhao2024galore} projects optimiser state, not adapter capacity,
onto a low-rank subspace; it tackles the memory cost of full-parameter
training rather than the allocation question we address. To our
knowledge no published LoRA variant uses an explicit per-rank parameter-
or activation-cost denominator in its importance score.

\section{Method}
\label{sec:method}

\subsection{Allocator rule}
\label{sec:method-rule}

Let $\theta_i = (A_i, B_i)$ be the LoRA parameters at target module $i$,
with $A_i \in \mathbb{R}^{r \times d_i^{\text{in}}}$ and
$B_i \in \mathbb{R}^{d_i^{\text{out}} \times r}$. During warmup we
maintain an EMA over the Frobenius gradient norm:
\begin{equation}
g_i^{(t)} \;=\; \beta\, g_i^{(t-1)} \;+\; (1-\beta)
   \bigl(\|\nabla_{A_i} \mathcal{L}\|_F + \|\nabla_{B_i} \mathcal{L}\|_F\bigr),
\end{equation}
seeded at the first observation rather than at zero. The per-rank
parameter cost
\begin{equation}
c_i \;=\; d_i^{\text{in}} + d_i^{\text{out}}
\end{equation}
is exact for a single rank added to module $i$. The score
$s_i = g_i / c_i^{\alpha}$ is then mapped to integer ranks by a
budget-preserving allocator that floors at $r_{\min}$, distributes the
remaining budget proportionally to $s_i$, caps at $r_{\max}$, and
deterministically rebalances $\pm 1$ per step until $\sum_i r_i = B$
exactly.

\subsection{Two-stage training}
\label{sec:method-twostage}

Stage 1 trains a uniform-rank LoRA model for $W$ warmup steps,
accumulating $g_i$ via the hook above. At the warmup-to-stage-2
transition the allocator runs once inside a context manager that
charges its wall-clock to a \texttt{scheduler\_overhead\_seconds}
field in the hardware logger. The warmup LoRA weights are then
\emph{discarded}: stage 2 rebuilds the LoRA model from a fresh base
with PEFT's \texttt{LoraConfig.rank\_pattern} set from the allocator's
output, then trains to completion on the same hyperparameters as
stage 1. The intent of warmup is to surface gradient signal, not to
pre-train weights.

A dynamic in-training variant --- re-running the allocator periodically
rather than once --- is out of scope. Resizing rank mid-training would
require reshaping the optimiser state of every adapter and is left to
future work.

\begin{algorithm}[t]
\caption{Hardware-Aware Two-Stage LoRA}
\label{alg:hwlora}
\begin{algorithmic}
\STATE \textbf{Input:} base model, dataloader, total budget $B$,
       $r_{\min}$, $r_{\max}$, $\alpha$, EMA decay $\beta$,
       warmup steps $W$, total steps $T$
\STATE Build LoRA model with uniform rank $B/n$ over $n$ targets
\STATE Initialise $g_i \leftarrow \perp$ for each target
\FOR{$t = 1$ to $W$}
   \STATE Forward, backward, update $g_i$ via Frobenius EMA
   \STATE Optimiser step
\ENDFOR
\STATE With \texttt{scheduler\_block} timing context:
\STATE \quad Compute $s_i = g_i / c_i^{\alpha}$
\STATE \quad Allocate $\{r_i\}$ s.t.\ $\sum_i r_i = B$,
       $r_{\min} \le r_i \le r_{\max}$
\STATE Reload base model; build LoRA with \texttt{rank\_pattern}
       $= \{r_i\}$
\FOR{$t = W{+}1$ to $T$}
   \STATE Forward, backward, optimiser step
\ENDFOR
\end{algorithmic}
\end{algorithm}

\section{Experimental Setup}
\label{sec:setup}

\textbf{Task.} SST-2 binary sentiment classification (67k train / 872
val) from GLUE \cite{wang2018glue, socher2013sst}.

\textbf{Model.} \texttt{distilbert-base-uncased}
\cite{sanh2019distilbert}.

\textbf{LoRA targets.} Four projection types --- \texttt{q\_lin},
\texttt{v\_lin}, \texttt{lin1}, \texttt{lin2} --- giving 24 modules
total: 12 attention projections at parameter cost
$c_i = 768 + 768 = 1536$ and 12 FFN projections at
$c_i = 768 + 3072 = 3840$. The $2.5\times$ cost heterogeneity is what
makes the gradient-only ablation non-degenerate; on attention-only
targets every $c_i$ is identical, and $\alpha{=}0$ vs.\ $\alpha{=}1$
produce bit-identical allocations.

\textbf{Rank budget.} $B = 192$ (24 modules $\times$ uniform rank 8),
invariant across all four methods so that any difference is not
attributable to using more trainable parameters.

\textbf{Methods.} Uniform LoRA, AdaLoRA, Gradient-Adaptive LoRA
($\alpha{=}0$), Hardware-Aware LoRA ($\alpha{=}1$). The $\alpha$-sweep
ablation additionally includes $\alpha{=}0.5$.

\textbf{Shared.} Identical optimiser (AdamW), learning rate
($2 \times 10^{-4}$), batch size (32), maximum steps (3 epochs), and
training loop for all methods, enforced by a single dispatcher in
\texttt{src/train.py}. AdaLoRA's per-step \texttt{update\_and\_allocate}
is plumbed through the same loop's \texttt{post\_step\_hook} so its
overhead is also charged to \texttt{scheduler\_overhead\_seconds}.

\textbf{Seeds.} 42, 43, 44.

\textbf{Hardware.} Single NVIDIA RTX 4070 Laptop GPU, CUDA 12.4,
PyTorch 2.6.0+cu124. The full sweep ran in $15780$ s ($\sim$$4.4$ h)
over $5 \times 3 = 15$ runs, all successful. Aggregation deduplicates by
$(\text{method}, \text{seed}, \alpha)$, keeping the lexicographically
latest \texttt{run\_id}. Four pre-provenance logs missing the
self-describing config row are skipped and contribute to no reported
number.

\section{Results}
\label{sec:results}

\subsection{Statistical performance}
\label{sec:results-stat}

\begin{table}[t]
\caption{Final validation loss, accuracy, and steps to reach the 0.90
accuracy target on SST-2. Mean $\pm$ std over three seeds. Parenthesised
counts are the number of seeds (out of 3) that crossed the target.}
\label{tab:stat}
\vskip 0.05in
\centering
\small
\begin{tabular}{lccc}
\toprule
Method & Final Val Loss & Final Val Acc.\ & Steps to 0.90 \\
\midrule
Uniform LoRA              & $0.286 \pm 0.016$ & $0.902 \pm 0.006$ & $1400 \pm 300$ (3/3) \\
AdaLoRA                   & $0.268 \pm 0.003$ & $0.894 \pm 0.003$ & never (0/3) \\
Gradient-Adaptive ($\alpha{=}0$) & $0.287 \pm 0.008$ & $0.904 \pm 0.005$ & $1333 \pm 231$ (3/3) \\
Hardware-Aware ($\alpha{=}1$)    & $0.282 \pm 0.007$ & $\mathbf{0.906 \pm 0.001}$ & $1767 \pm 351$ (3/3) \\
\bottomrule
\end{tabular}
\vskip -0.05in
\end{table}

\begin{figure}[t]
\centering
\includegraphics[width=\linewidth]{../results/figures/val_accuracy_vs_walltime.png}
\caption{Validation accuracy versus training wall-clock for each method.
Mean over three seeds; the four methods land in a tight band but
arrive at different rates. Generated by \texttt{src/metrics.py}.}
\label{fig:acc-walltime}
\end{figure}

Table~\ref{tab:stat} reports the four methods' final statistical
performance and Fig.~\ref{fig:acc-walltime} their accuracy trajectories
against wall-clock. All four land in a tight $0.894$--$0.906$ accuracy
band; on SST-2 with DistilBERT the differences are within seed variance.
Hardware-Aware shows the lowest seed-to-seed variance ($\pm 0.001$ over
three seeds) and the highest mean accuracy. AdaLoRA produces the lowest
validation loss ($0.268$) but does not cross the $0.90$ accuracy
threshold in any of three seeds within three epochs, while the other
three methods cross it in all three. We report this factually without
further causal claim.

\subsection{Hardware performance}
\label{sec:results-hw}

\begin{table*}[t]
\caption{Hardware metrics. Peak memory is
\texttt{torch.cuda.max\_memory\_allocated}; throughput is the per-method
EMA over training. Mean $\pm$ std over three seeds. Trainable parameters
is deterministic per-method given the rank pattern.}
\label{tab:hw}
\vskip 0.05in
\centering
\small
\begin{tabular}{lcccc}
\toprule
Method & Peak Memory (MB) & Examples/sec & Wall-Clock (s) & Trainable Params \\
\midrule
Uniform LoRA              & $\mathbf{948.8 \pm 1.0}$  & $920.3 \pm 15.4$ & $1011.4 \pm 1.5$ & $1{,}108{,}226$ \\
AdaLoRA                   & $964.3 \pm 0.1$  & $601.7 \pm 25.9$ & $1130.2 \pm 1.0$ & $1{,}366{,}562$ \\
Gradient-Adaptive         & $1210.7 \pm 1.3$ & $918.7 \pm 18.4$ & $1024.1 \pm 1.7$ & $1{,}156{,}610$ \\
Hardware-Aware            & $1208.3 \pm 1.2$ & $\mathbf{938.4 \pm 44.5}$ & $1025.6 \pm 0.6$ & $\mathbf{1{,}081{,}346}$ \\
\bottomrule
\end{tabular}
\vskip -0.05in
\end{table*}

\begin{figure}[t]
\centering
\includegraphics[width=\linewidth]{../results/figures/peak_memory_bars.png}
\caption{Peak GPU memory by method. Two-stage methods (Hardware-Aware,
Gradient-Adaptive) carry $\sim$$250$\,MB more peak memory than Uniform
despite using fewer or comparable parameters; the gap is activation,
not parameter, memory.}
\label{fig:peakmem}
\end{figure}

\begin{figure}[t]
\centering
\includegraphics[width=\linewidth]{../results/figures/examples_per_second_bars.png}
\caption{Training throughput by method. AdaLoRA's per-step
\texttt{update\_and\_allocate} hook costs $\sim$$35\%$ of throughput;
the three single-allocator methods cluster around $\sim$$920$
examples/s.}
\label{fig:eps}
\end{figure}

At the same $192$-rank budget, Hardware-Aware uses
$\mathbf{1{,}081{,}346}$ trainable parameters --- fewer than Uniform
($1{,}108{,}226$), Gradient-Adaptive ($1{,}156{,}610$), and AdaLoRA
($1{,}366{,}562$). This is the cost penalty doing what it is designed
to do: rank concentrates on the cheaper attention modules
($c_i{=}1536$) rather than the more expensive FFN modules
($c_i{=}3840$). Throughput splits cleanly along scheduler design
(Fig.~\ref{fig:eps}): Uniform, Gradient-Adaptive, and Hardware-Aware
all run at $\sim$$920$ examples/s, while AdaLoRA runs at $602$~ex/s ---
a $\sim$$35\%$ throughput penalty attributable to its per-step
\texttt{update\_and\_allocate} work.

Peak memory tells the opposite story to the one we expected from the
parameter counts (Fig.~\ref{fig:peakmem}). Uniform sits at $949$\,MB
and AdaLoRA at $964$\,MB, but Hardware-Aware reaches $1208$\,MB and
Gradient-Adaptive $1211$\,MB. Both two-stage methods carry
$\sim$$250$\,MB more peak memory than Uniform despite Hardware-Aware
using fewer parameters. \texttt{torch.cuda.max\_memory\_allocated}
reflects activation and optimiser-state peaks during training, not
parameter count alone; rank concentrating on attention grows those
layers' $A$/$B$ factors and their activation footprint along the
forward/backward pass. Our cost function $c_i = d_i^{\text{in}} +
d_i^{\text{out}}$ does not currently model activation memory --- a
direction for follow-up rather than a closed story.

\subsection{Systems tradeoff}
\label{sec:results-systems}

\begin{table}[t]
\caption{Systems tradeoff. Accuracy per MB of peak memory, time to
reach the 0.90 accuracy target, and scheduler overhead (the wall-clock
attributable to allocation work, charged via the logger's
\texttt{scheduler\_block} context).}
\label{tab:systems}
\vskip 0.05in
\centering
\small
\begin{tabular}{lccc}
\toprule
Method & Acc.\ per MB & Time to 0.90 (s) & Sched.\ Overhead (s) \\
\midrule
Uniform LoRA   & $\mathbf{0.00095 \pm 0.00001}$ & $\mathbf{225.2 \pm 45.3}$ (3/3) & $0.00 \pm 0.00$ \\
AdaLoRA        & $0.00093 \pm 0.00000$ & never (0/3) & $76.67 \pm 0.33$ \\
Hardware-Aware & $0.00075 \pm 0.00000$ & $293.5 \pm 54.1$ (3/3) & $\mathbf{3.98 \pm 0.13}$ \\
\bottomrule
\end{tabular}
\vskip -0.05in
\end{table}

\begin{figure}[t]
\centering
\includegraphics[width=\linewidth]{../results/figures/scheduler_overhead_bars.png}
\caption{Scheduler overhead by method. Hardware-Aware's one-shot
allocator runs at the warmup-to-stage-2 transition only; AdaLoRA's
\texttt{update\_and\_allocate} runs every training step. The gap is
structural and would widen with longer training.}
\label{fig:sched}
\end{figure}

The scheduler-overhead column of Table~\ref{tab:systems} is the
project's clearest systems-flavour finding. Hardware-Aware's reallocator
runs once at the warmup-to-stage-2 transition and costs
$\mathbf{4.0 \pm 0.13}$\,s of wall-clock, while AdaLoRA's runs every
training step and accumulates to $\mathbf{76.7 \pm 0.33}$\,s ---
roughly $19\times$ more (Fig.~\ref{fig:sched}). This gap is structural
to one-shot vs.\ continuous reallocation and would only widen on longer
training horizons. Time-to-target is mixed: Uniform reaches $0.90$ in
$225$\,s vs.\ Hardware-Aware's $294$\,s, and AdaLoRA does not reach it
at all in three seeds. Accuracy-per-MB favours Uniform ($0.00095$)
over Hardware-Aware ($0.00075$) --- a direct consequence of the
peak-memory gap above, not a property of the allocator itself.

We follow CLAUDE.md's table convention: this systems tradeoff table
limits to the three method-vs-method baselines (Uniform, AdaLoRA,
Hardware-Aware). Gradient-Adaptive is an ablation rather than a
comparator, and is reported in Tables~\ref{tab:stat}, \ref{tab:hw} and
in the $\alpha$-sweep below.

\subsection{$\alpha$-sweep ablation}
\label{sec:results-alpha}

\begin{table}[t]
\caption{$\alpha$-sweep at fixed budget and gradient signal. The
``Attention rank share'' column is the fraction of total rank assigned
to the 12 attention modules ($c_i{=}1536$) versus the 12 FFN modules
($c_i{=}3840$). Mean $\pm$ std over three seeds.}
\label{tab:alpha}
\vskip 0.05in
\centering
\small
\begin{tabular}{lccc}
\toprule
$\alpha$ & Final Val Acc.\ & Time to 0.90 (s) & Attn.\ Rank Share \\
\midrule
$0.0$ (gradient-only)        & $0.904 \pm 0.005$ & $227.6 \pm 35.9$ (3/3) & $0.391 \pm 0.005$ \\
$0.5$                        & $0.906 \pm 0.003$ & $320.1 \pm 85.1$ (3/3) & $0.476 \pm 0.008$ \\
$1.0$ (full hardware penalty)& $0.906 \pm 0.001$ & $293.5 \pm 54.1$ (3/3) & $0.561 \pm 0.003$ \\
\bottomrule
\end{tabular}
\vskip -0.05in
\end{table}

\begin{figure*}[t]
\centering
\includegraphics[width=0.95\linewidth]{../results/figures/rank_allocation_heatmap.png}
\caption{Per-module rank allocation across methods, averaged over three
seeds. Rows: the four target projection types
(\texttt{attn.q\_lin}, \texttt{attn.v\_lin}, \texttt{ffn.lin1},
\texttt{ffn.lin2}), aggregated across the six DistilBERT layers.
Columns: the four methods plus the $\alpha$-sweep variants. The
attention-rank-share growth in Table~\ref{tab:alpha} is driven primarily
by \texttt{attn.v\_lin}; \texttt{attn.q\_lin} stays low across every
$\alpha$ tested.}
\label{fig:heatmap}
\end{figure*}

The $\alpha$-sweep is the cleanest single piece of evidence that the
cost term in $s_i = g_i / c_i^{\alpha}$ is doing real work
(Table~\ref{tab:alpha}). As $\alpha$ moves from $0.0$ to $0.5$ to $1.0$,
attention rank share grows monotonically from
$\mathbf{39.1\% \to 47.6\% \to 56.1\%}$ while the gradient signal is
held fixed across the three configurations. Final accuracy stays
essentially flat across the sweep ($0.904$, $0.906$, $0.906$), so the
systems reshaping is effectively free on accuracy for this task.

\subsection{Asymmetric skew across attention projections}
\label{sec:results-skew}

The heatmap (Fig.~\ref{fig:heatmap}) reveals that the attention-share
growth above is asymmetric across the two attention projections: rank
migrates from \texttt{ffn.lin1} to \texttt{attn.v\_lin} as $\alpha$
grows, while \texttt{attn.q\_lin} stays low across all methods (its
gradient signal is too small for the cost penalty to lift it) and
\texttt{ffn.lin2} stays low throughout. The $39\% \to 56\%$
attention-share growth in Table~\ref{tab:alpha} is driven primarily by
\texttt{v\_lin}.

This exposes a structural property of the rule
$s_i = g_i / c_i^{\alpha}$: the cost term reshapes allocations
\emph{among} modules with non-trivial gradient signal, but it cannot
manufacture importance for modules whose $g_i$ is near zero. The
penalty is multiplicative, not additive. Three implications follow:

\textbf{1. ``Hardware-aware'' is more precisely ``hardware-biased
gradient'' allocation.} The cost term shapes the distribution of rank
among modules that already have learning signal, but does not recruit
modules from a wider candidate pool. This is a sharper claim than
``the cost term reshapes allocations'' and is the one supported by the
data.

\textbf{2. It likely explains the small accuracy gap between the two
two-stage variants.} Hardware-Aware ($\alpha{=}1$) and Gradient-Adaptive
($\alpha{=}0$) end at $0.904$--$0.906$ because they allocate among the
same set of gradient-signal-rich modules and differ only in the
proportions. If the cost penalty reached \texttt{attn.q\_lin}, we would
expect a larger accuracy delta in either direction; its absence is
consistent with redistribution-within-shortlist rather than
expansion-of-shortlist.

\textbf{3. It suggests an additive variant.} A rule of the form
$s_i = g_i + \lambda / c_i$ would let the cost term contribute
independently of $g_i$ and could route rank to low-gradient modules.
Whether that is desirable is a separate question --- funding modules
that are not learning may waste budget --- but it would disambiguate
whether \texttt{attn.q\_lin}'s low rank reflects genuine
uninformativeness or insufficient warmup signal.

\section{Limitations}
\label{sec:limitations}

\begin{itemize}
\item \textbf{Single task and single model.} SST-2 with DistilBERT-base
      only. Whether the allocator's behaviour generalises to multi-class
      GLUE tasks or to RoBERTa- or Llama-class models is not validated
      here.
\item \textbf{Two-stage one-shot reallocation only.} A dynamic
      in-training variant is deferred. We do not know whether
      continuous reallocation would close the time-to-target gap with
      Uniform LoRA on this task.
\item \textbf{Three seeds.} Too few to support paired significance
      testing. We report mean $\pm$ standard deviation and let the
      reader judge.
\item \textbf{Mixed practical-efficiency picture vs.\ Uniform.}
      Hardware-Aware uses fewer trainable parameters and achieves the
      lowest seed variance, but Uniform LoRA reaches the $0.90$ target
      faster and uses less peak memory on this task. The headline
      systems wins (throughput, scheduler overhead) are observed
      against AdaLoRA, not against Uniform.
\item \textbf{Peak memory not modeled in the cost function.}
      $c_i = d_i^{\text{in}} + d_i^{\text{out}}$ captures parameter
      cost per rank, not activation memory. The unexpected peak-memory
      excess of the two-stage methods over Uniform is an open
      follow-up.
\item \textbf{Multiplicative penalty cannot rescue near-zero gradient
      signal.} As \S\ref{sec:results-skew} shows, \texttt{attn.q\_lin}
      stays low-rank across every $\alpha$ we tested because its $g_i$
      is small; no value of $\alpha$ can route rank to a module with
      no learning signal. An additive variant would behave differently.
\end{itemize}

\section{Conclusion}
\label{sec:conclusion}

The cost term in $s_i = g_i / c_i^{\alpha}$ redistributes rank
meaningfully among gradient-signal-rich modules --- the $\alpha$-sweep
monotonicity demonstrates this independent of any change to the
gradient signal itself --- and yields the lowest trainable-parameter
count of any method evaluated, competitive accuracy with the lowest
seed variance, and a $\sim$$19\times$ reduction in scheduler overhead
relative to AdaLoRA. It does not recruit rank for modules with
near-zero gradient signal (the rule is multiplicative), and it does
not deliver an across-the-board practical-efficiency win over Uniform
LoRA on this task, with peak memory and time-to-target as the headline
tradeoffs.

Our method explores whether hardware-aware rank allocation can improve
practical efficiency compared with AdaLoRA under constrained training
budgets. The evidence presented here supports a sharpened version of
that claim: the cost term shifts \emph{where} rank goes among modules
already contributing gradient signal, at essentially zero accuracy cost
on SST-2 with DistilBERT. Whether that reshaping translates to
practical wins at larger scale, on tasks with richer gradient signal
across more modules, or under additive cost rules, is left to future
work.

\bibliography{references}
\bibliographystyle{icml2025}

\end{document}
